% !PRINTED 146
\weberpage{146}
%<<CONT
the same subsidiary group, say in $Q\pi_1$, and consequently there are two
permutations $\varkappa_1, \varkappa_2$ in $Q$ which satisfy the condition

\[ \varkappa_1 \pi_1 = (0, 1), \quad \varkappa_2 \pi_1 = (0, 2) \]

It is therefore

\[ \varkappa_1 \pi_1 \pi_1^{-1} \varkappa_2^{-1} = \varkappa_1 \varkappa_2^{-1} = (0,1)(0,2) = (0,1,2) \]

contained in $Q$. But this contradicts our assumption that $Q$ is to contain no
three-term cycle. Accordingly $j$ cannot be $< n$. But if $j = n$, and hence the
degree of $Q$ equal to $\Pi(n-1)$, then the $n-1$ permutations (5) must occur in
$n-1$ different subsidiary groups, and $P$ can be represented thus:

\weq{6}{P = Q + Q(0,1) + Q(0,2) + \cdots + Q(0, n-1).}

If we consider any other transposition, e.g. $(2,3)$, this cannot occur
in $Q$, nor in $Q(0,2)$ or in $Q(0,3)$, since otherwise
$(2,3)(0,2) = (0,2,3)$ or $(2,3)(0,3) = (0,3,2)$ would occur in $Q$. We may therefore,
without impairing the generality, assume that $(2,3)$ occurs in $Q(0,1)$,
and from this it follows:

\gesperrt{$\beta$) The group $Q$ contains the transposition pair}

\weq{7}{(0,1)(2,3).}

The group $Q$ has a divisor $Q_0$, which consists of all those permutations
that leave the figure 0 in its place. Since $Q$ is transitive, it follows that if
$\varkappa_1, \varkappa_2, \ldots, \varkappa_{n-1}$ are elements of $Q$ which carry the 0
into 1, into 2, \ldots, into $n-1$, then

\weq{8}{Q = Q_0 + Q_0 \varkappa_1 + \cdots + Q_0 \varkappa_{n-1},}

and since $Q$ is of degree $\Pi(n-1)$, it follows from this that $Q_0$ is of degree

\weq{9}{g = \frac{\Pi(n-1)}{n}}

Since $g$ must be an integer, and hence $\Pi(n-1)$ divisible by $n$, we
conclude first of all that the case $j = n$ can never occur when $n$
is a prime, and for this case our Theorem II. is therefore proved.

In general, however, we can conclude:

\gesperrt{$\gamma$) If $n$ is not $= 4$, then the $n-1$ figures
$1, 2, \ldots, n-1$ are still transitively connected by $Q_0$.}
